\documentclass[11pt]{article}

% --- Math ---
\usepackage{amsmath, amssymb, amsthm, amsfonts}
\usepackage{mathrsfs}

\usepackage[utf8]{inputenc}
\usepackage[T1]{fontenc}
\usepackage{siunitx}

% --- Algorithms ---
\usepackage{algorithm}
\usepackage{algpseudocode}

% --- Tables ---
\usepackage{tabularx}
\usepackage{array}
\usepackage{booktabs}
\usepackage{multirow}
\usepackage{makecell}
\usepackage{ragged2e}
\newcolumntype{Y}{>{\RaggedRight\arraybackslash}X}

% --- Figures ---
\usepackage{graphicx}
\usepackage{float}

% --- Bibliography ---
\usepackage[authoryear]{natbib}

\pagestyle{plain}
\setlength{\textwidth}{6.5in}
\setlength{\oddsidemargin}{0in}
\setlength{\evensidemargin}{0in}
\setlength{\textheight}{9.2in}
\setlength{\topmargin}{-0.2in}
\setlength{\headheight}{0in}
\setlength{\headsep}{0in}
\setlength{\footskip}{.4in}
\renewcommand{\baselinestretch}{1.5}

\newtheorem{theorem}{Theorem}
\newtheorem{obs}{Observation}
\newtheorem{proposition}{Proposition}
\newtheorem{remark}{Remark}
\newtheorem{definition}{Definition}

% Compile safely even when the chamber figure is not in the same folder.
\newcommand{\safechamberfigure}{%
  \IfFileExists{chamber_structure.png}{%
    \includegraphics[width=0.88\textwidth]{chamber_structure.png}%
  }{%
    \fbox{\begin{minipage}{0.78\textwidth}
    \centering
    \vspace{8mm}
    Chamber structure figure placeholder.\\
    Place \texttt{chamber\_structure.png} in the same folder to display the figure.
    \vspace{8mm}
    \end{minipage}}%
  }%
}

\begin{document}

\begin{center}
{\Large \bf Integrated Sample Size Optimization and Test Chamber Scheduling under Environmental Constraints}\\
\vskip 12pt
\normalsize
Nur Per\c{c}in$^{1}$,
Ali Ekici$^{1}$\\[6pt]
\footnotesize
$^{1}$Department of Industrial Engineering, Ozyegin University, Istanbul, Turkiye\\[3pt]
\texttt{nur.percin@ozu.edu.tr, ali.ekici@ozyegin.edu.tr}
\end{center}

\begin{abstract}
Research and development testing of refrigerators requires the coordinated use of heterogeneous test chambers under temperature, humidity, voltage, stage-precedence, and due-date constraints. In current practice, sample sizes are often fixed, although sample decisions directly determine both workload and opportunities for parallel execution. This study introduces the Test Chamber Scheduling with Sample Optimization Problem (TCSSOP), in which product sample sizes, chamber-environment assignments, and resource-feasible schedules are optimized jointly.

An iterative two-stage framework is proposed. For a given sample vector, chamber environments are assigned by balancing environment-specific workload and voltage-sensitive workload over compatible chamber capacity. Given chamber settings, sample sizes are improved by local search and evaluated through an earliest-due-date constructive scheduler with explicit feasibility validation. When the search stagnates, a Variable Neighborhood Search (VNS) mechanism perturbs sample sizes using fixed fallback ratios and re-optimizes the resulting workload structure.

Computational experiments on 90 benchmark scenarios show that endogenous sample optimization substantially reduces total tardiness relative to the fixed three-sample policy. The best configuration, using a \(0.20\)--\(0.30\) VNS fallback policy with \(K_{\max}=200\), provides the strongest overall performance. The results show that sample allocation should be treated as a tactical planning decision and coordinated with chamber reconfiguration rather than fixed uniformly across projects.
\end{abstract}

\textbf{Keywords:} Test chamber scheduling; Sample size optimization; Decision-dependent workload; Resource-constrained scheduling; Variable neighborhood search; Multi-stage testing systems

\section{Introduction}
\label{sec:intro}

Refrigerator manufacturers operate under compressed product-development cycles, strict market-specific compliance requirements, and high pressure to launch new models on time. Before commercialization, products must complete design approval tests covering performance, safety, energy efficiency, and consumer-usage conditions \citep{IEA2021,IEC60335-1,IEC60335-2-24,IEC62552}. Delays in these tests postpone certification and production release, while inefficient use of chambers lowers the return on capital-intensive R\&D infrastructure \citep{OECD_MSTI2024,OECD_MSTI2025}.

The testing system considered in this study is motivated by refrigerator R\&D operations at Vestel. Products compete for a set of heterogeneous climate chambers, each containing several stations that operate under a common environmental configuration. Test operations require specific temperature and humidity conditions; some products also require voltage-compatible execution. In addition, the testing process follows stage-state logic: gas amount determination precedes pull-down testing, downstream tests are released after the pull-down phase, and final consumer-usage tests require the relevant downstream-stage release conditions to be satisfied.

A key operational feature is the use of physical samples. Industrial practice commonly applies a fixed three-sample policy. This rule is simple, but it ignores the fact that sample size is itself a scheduling decision. Pull-down tests scale with the number of samples and therefore increase upstream workload, whereas additional samples can enable downstream tests to run in parallel and reduce completion times. The effect of sample size on tardiness is consequently non-monotone and depends on due dates, chamber eligibility, environmental bottlenecks, and stage-state interactions.

We define the resulting problem as the \emph{Test Chamber Scheduling with Sample Optimization Problem} (TCSSOP). TCSSOP combines heterogeneous parallel resources, chamber-level environmental assignment, humidity and voltage eligibility, multi-stage test logic, and decision-dependent workload generation. The objective is to minimize total tardiness, with total sample usage used as a secondary lexicographic criterion.

The proposed solution framework is designed around the feedback between sample decisions and chamber configuration. A change in sample sizes modifies the environment-specific workload distribution, especially through pull-down operations. This change may make the previous chamber assignment unsuitable. Conversely, a different chamber assignment may change which products benefit most from additional samples. For this reason, the framework alternates between chamber assignment and sample-size optimization rather than solving them once in isolation.

This study makes three contributions. First, it benchmarks the fixed three-sample policy against endogenous sample-size optimization. Second, it develops an iterative framework that coordinates sample-size decisions with chamber-environment reassignment and feasible scheduling. Third, it introduces a VNS diversification strategy that perturbs sample-size decisions rather than only sequencing or assignment decisions. The computational results show that these components jointly produce substantial tardiness reductions across 90 benchmark scenarios.

\section{Literature Review}
\label{sec:literature_review}

The problem is related to several scheduling and test-planning streams. Resource-constrained multi-project scheduling and unrelated parallel-machine scheduling provide the closest scheduling background because multiple projects compete for heterogeneous resources under precedence, due-date, and eligibility restrictions. Recent surveys emphasize that practical scheduling systems increasingly require integrated treatment of resource heterogeneity, precedence, and operational constraints \citep{HartmannBriskorn2022,GomezSanchezLallaRuizFernandezGilCastroVoss2023}. In unrelated parallel-machine environments, machine eligibility and tardiness-oriented objectives substantially increase computational difficulty and motivate specialized heuristics \citep{UlagaDurasevicJakobovic2022,MaeckerShenMonch2023}.

Resource-constrained multi-project scheduling provides a natural reference point because several projects are processed simultaneously and share limited resources. However, standard RCMPSP formulations usually assume that the set of activities and their processing requirements are known in advance. In TCSSOP, the activity set is partly induced by the sample vector. Pull-down operations are generated per sample, and downstream parallelism depends on the number of available physical units. This workload-generation feature is not captured by conventional multi-project scheduling models.

Industrial scheduling often also involves setup or configuration decisions. The survey by \citet{Allahverdi2015} shows that ignoring setup-related effects can produce unrealistic schedules. In test-chamber systems, chamber environmental settings play a similar role because each chamber can serve only tests compatible with its current temperature-humidity configuration. This creates a reconfigurable resource-allocation problem related in spirit to reconfigurable manufacturing systems \citep{Koren1999,MehrabiUlsoyKoren2000}, although the present setting concerns testing rather than production.

Integrated decision-making models such as order acceptance and scheduling combine workload selection with sequencing under capacity limits \citep{Slotnick2011,WangYe2019}. Reliability demonstration testing optimizes sample size and test duration under statistical risk or cost criteria \citep{Fernandez2022}, but typically abstracts away detailed chamber capacity and stage-dependent scheduling. TCSSOP lies between these areas: sample sizes generate the workload, while the resulting operations must be scheduled on constrained heterogeneous chambers.

Methodologically, VNS and related neighborhood-based metaheuristics have performed well in complex scheduling settings because they combine local improvement with systematic diversification \citep{MejiaYuraszeck2020}. The diversification mechanism proposed here differs from conventional VNS designs by perturbing sample-size decisions. Thus, the search explores alternative workload structures, which then induce different chamber-utilization and scheduling regimes.

\section{Problem Definition}
\label{sec:pd}

\subsection{System description}
\label{subsec:system}

The facility contains a set of chambers \(\mathcal{C}\). Chamber \(c\) has \(M_c\) stations and supports temperature adjustment; only selected chambers support humidity adjustment \(h_c\) and voltage adjustment \(v_c\). Stations in the same chamber operate under one common environment \(X_c\), where an environment represents a temperature-humidity combination.

The set of products is denoted by \(\mathcal{P}\). Each product \(p\) has a due date \(d_p\), a voltage requirement indicator \(\nu_p\), and a product-specific required test set \(\mathcal{T}_p \subseteq \mathcal{T}\). Test requirements are determined by target market, climatic class, energy-efficiency regulation, optional engineering decisions, and electrical specification. Table~\ref{tab:test_attributes_short} summarizes the test catalog used in the study.

\begin{table}[H]
\caption{List of tests and their environmental conditions}
\label{tab:test_attributes_short}
\centering
\footnotesize
\setlength{\tabcolsep}{3.5pt}
\begin{tabular}{|c|c|l|c|c|c|c|}
\hline
\textbf{Stage} & \textbf{Test No.} & \textbf{Test Name} & \textbf{Temperature} & \textbf{Ambient Humidity} & \textbf{Samples} & \textbf{Duration (days)} \\ \hline
1 & 1 & Gas Amount Determination & 43/38/32~\si{\degreeCelsius} & Ambient & 1 & 10 \\ \hline
2 & 2  & Pull-down & 43/38/32~\si{\degreeCelsius} & Ambient & All samples & 5 \\ \hline
3 & 3  & Energy Efficiency & \SI{16}{\degreeCelsius} & Ambient & 1 & 15 \\ \hline
3 & 4  & Energy Efficiency & \SI{25}{\degreeCelsius} & Ambient & 1 & 15 \\ \hline
3 & 5  & Energy Efficiency & \SI{32}{\degreeCelsius} & Ambient & 1 & 15 \\ \hline
3 & 6$^*$ & Performance & \SI{10}{\degreeCelsius} & Ambient & 1 & 10 \\ \hline
3 & 7$^*$ & Performance & \SI{16}{\degreeCelsius} & Ambient & 1 & 10 \\ \hline
3 & 8  & Performance & \SI{25}{\degreeCelsius} & Ambient & 1 & 10 \\ \hline
3 & 9$^{**}$ & Performance & \SI{32}{\degreeCelsius} & Ambient & 1 & 10 \\ \hline
3 & 10$^{**}$ & Performance & \SI{38}{\degreeCelsius} & Ambient & 1 & 10 \\ \hline
3 & 11$^{**}$ & Performance & \SI{43}{\degreeCelsius} & Ambient & 1 & 10 \\ \hline
4 & 12$^{***}$ & Freezing Capacity & \SI{25}{\degreeCelsius} & Ambient & 1 & 10 \\ \hline
4 & 13$^{***}$ & Temperature Rise & \SI{25}{\degreeCelsius} & Ambient & 1 & 5 \\ \hline
5 & 14 & Consumer Usage & \SI{10}{\degreeCelsius} & Ambient & 1 & 15 \\ \hline
5 & 15 & Consumer Usage & \SI{25}{\degreeCelsius} & 85\% RH & 1 & 15 \\ \hline
5 & 16 & Consumer Usage & \SI{32}{\degreeCelsius} & 85\% RH & 1 & 15 \\ \hline
\end{tabular}
\vspace{1mm}
\begin{minipage}{0.95\textwidth}
\footnotesize
$^*$ One of these tests is selected based on target market climatic class requirements.\\
$^{**}$ One of these tests is selected based on target market climatic class requirements.\\
$^{***}$ Optional tests decided by the product engineer.
\end{minipage}
\end{table}

The chamber system includes 28 chambers and 240 stations. Chamber capacities range from 6 to 12 stations. Humidity and voltage capabilities are heterogeneous, so the effective capacity available for a given environment may be much smaller than the total station count. Figure~\ref{fig:chamber_3d_short} illustrates the chamber-station structure.

\begin{figure}[H]
\centering
\safechamberfigure
\caption{Test chambers consist of multiple stations operating under a shared chamber-level environment.}
\label{fig:chamber_3d_short}
\end{figure}

\begin{table}[H]
\centering
\caption{Chamber capacities and technical adjustment capabilities}
\label{tab:chamber_capabilities_short}
\footnotesize
\setlength{\tabcolsep}{4pt}
\renewcommand{\arraystretch}{1.1}
\begin{tabular}{lcccc}
\hline
\textbf{Chamber} &
\shortstack{\textbf{No. of}\\\textbf{Stations}} &
\shortstack{\textbf{Temperature}\\\textbf{Adjustment}} &
\shortstack{\textbf{Humidity}\\\textbf{Adjustment}} &
\shortstack{\textbf{Voltage}\\\textbf{Adjustment}} \\
\hline
T001 & 8  & 1 & 0 & 0 \\
T002 & 8  & 1 & 1 & 1 \\
T003 & 8  & 1 & 1 & 1 \\
T004 & 8  & 1 & 1 & 1 \\
T005 & 8  & 1 & 1 & 1 \\
T006 & 8  & 1 & 0 & 0 \\
T007 & 8  & 1 & 0 & 0 \\
T008 & 12 & 1 & 0 & 1 \\
T009 & 12 & 1 & 0 & 0 \\
T010 & 12 & 1 & 1 & 0 \\
T011 & 12 & 1 & 1 & 0 \\
T012 & 12 & 1 & 1 & 0 \\
T013 & 12 & 1 & 1 & 0 \\
T014 & 6  & 1 & 1 & 0 \\
T015 & 6  & 1 & 1 & 0 \\
T016 & 6  & 1 & 0 & 1 \\
T017 & 6  & 1 & 0 & 1 \\
T018 & 6  & 1 & 0 & 1 \\
T019 & 6  & 1 & 0 & 1 \\
T020 & 6  & 1 & 0 & 1 \\
T021 & 6  & 1 & 0 & 1 \\
T022 & 6  & 1 & 0 & 1 \\
T023 & 6  & 1 & 0 & 1 \\
T024 & 8  & 1 & 1 & 0 \\
T025 & 8  & 1 & 1 & 1 \\
T026 & 12 & 1 & 0 & 1 \\
T027 & 12 & 1 & 0 & 1 \\
T028 & 12 & 1 & 0 & 1 \\
\hline
\end{tabular}
\end{table}

The chamber table highlights why TCSSOP cannot be reduced to a simple capacity-balancing problem. A workload peak at an ambient-temperature environment may be served by many chambers, but a workload peak involving humidity or voltage may be restricted to a much smaller subset. Therefore, the planning model must track both total station capacity and capability-specific station capacity.

\subsection{Test requirement generation}
\label{subsec:req_generation}

Product-specific test sets are generated from climatic class, energy regulation, optional engineering decisions, and voltage specification. Climatic class determines which performance-temperature alternatives are required. Energy-efficiency regulation determines whether a product requires a single-point or multi-point energy test. Optional tests such as freezing capacity and temperature rise are included only when specified by the product engineer.

The environmental requirement of a test is represented at the required test-operation level. In notation, \(e(t)\) denotes the environment of a selected test variant \(t\). When a product-specific requirement selects among alternative temperatures, the selected variant is treated as the corresponding test type in the catalog. Thus, the scheduler sees a binary required-test matrix over concrete test variants, each with a fixed environment.

Voltage is modeled at the product level rather than the test level because the electrical configuration follows the product's target market. If a product requires voltage-compatible execution, all of its operations must be processed on voltage-capable chambers. This one-sided condition reflects practice: a voltage-capable chamber can process a product without voltage requirements, but the reverse is not feasible.

\subsection{Decision-dependent workload}
\label{subsec:workload}

Each product has a sample-size decision \(s_p\), bounded by \(s_{\min} \le s_p \le s_{\max}\). Pull-down tests are executed once per sample, whereas other required tests are executed once per product. Thus, the multiplicity of test \(t\) for product \(p\) is
\[
m_{pt}(s_p)=
\begin{cases}
s_p, & \text{if }t\text{ is a pull-down test,}\\
1, & \text{otherwise.}
\end{cases}
\]
For environment \(e\), the induced workload is
\[
W_e(S)=
\sum_{p\in\mathcal{P}}
\sum_{\substack{t\in\mathcal{T}_p\\e(t)=e}}
m_{pt}(s_p)\tau_t,
\]
where \(\tau_t\) is the processing time of test \(t\). Voltage-sensitive workload \(W_e^V(S)\) is defined similarly by summing only products with \(\nu_p=1\).

This formulation captures the central trade-off. Increasing \(s_p\) increases pull-down workload but may reduce tardiness by enabling downstream parallelism. Reducing \(s_p\) saves workload but may serialize downstream testing on too few physical units.

\subsection{Stage-state logic and sample feasibility}
\label{subsec:stage_logic}

The scheduling logic is implemented through four job categories: gas amount determination, pull-down, other downstream tests, and consumer-usage tests. Gas amount determination is executed at most once for each product when it is required by the product's test matrix. If it is required, it is the first operation of the product and is released at time zero. If it is not required, the product is treated as already released for the pull-down/downstream logic.

Pull-down is executed once for each physical sample of a product. For product \(p\) with \(s_p\) samples, the pull-down phase therefore generates \(s_p\) sample-specific operations. Each pull-down operation can start only after the gas amount determination operation of the same product has finished. Different pull-down operations of the same product may run in parallel on different samples and compatible stations, subject to station availability and sample-level non-overlap.

The downstream non-consumer tests consist of the required energy-efficiency, performance, freezing-capacity, and temperature-rise tests. In the implemented configuration, these tests are released only after the full pull-down phase of the product has been completed, that is, after all sample-level pull-down operations have finished. Each required downstream non-consumer test is scheduled once per product and may be assigned to any available sample of that product after this release time.

Consumer-usage tests are modeled as the final category. A required consumer-usage test is released only after the pull-down phase has been completed and after all required downstream non-consumer tests of the same product have been started. Thus, the release time for consumer-usage tests is the maximum of the pull-down phase completion time and the latest start time among the required downstream non-consumer tests. Consumer-usage tests do not wait for those downstream non-consumer tests to finish.

Let \(G_p\) be the gas operation of product \(p\), \(PD_{pj}\) be the pull-down operation of sample \(j\), and \(\mathcal{O}_p\) and \(\mathcal{U}_p\) be the sets of required downstream non-consumer and consumer-usage tests, respectively. If gas and pull-down are required, then
\[
G_p \prec PD_{pj} \qquad \forall j=1,\ldots,s_p.
\]
The implemented downstream release rule is
\[
\max_{j=1,\ldots,s_p} C_{PD_{pj}} \le S_o
\qquad \forall o\in\mathcal{O}_p,
\]
where \(S_o\) denotes the start time of downstream non-consumer test \(o\). For each consumer-usage test \(u\in\mathcal{U}_p\),
\[
S_u \ge
\max\left\{
\max_{j=1,\ldots,s_p} C_{PD_{pj}},
\max_{o\in\mathcal{O}_p} S_o
\right\}.
\]
The second term is based on the start times of the downstream non-consumer tests, not their completion times. Sample feasibility is enforced independently from product-stage feasibility: at any time, a physical sample can process at most one operation.

\subsection{Feasibility conditions and objective}
\label{subsec:objective}

A feasible schedule must satisfy the following conditions:
\begin{enumerate}
    \item all and only required tests are scheduled with the multiplicities induced by \(S\);
    \item each operation is assigned to a chamber whose environment matches the test requirement;
    \item humidity-controlled tests are assigned only to chambers with \(h_c=1\);
    \item if \(\nu_p=1\), operations of product \(p\) are assigned only to chambers with \(v_c=1\);
    \item each station processes at most one operation at a time;
    \item each physical sample is used by at most one operation at a time;
    \item stage-state logic is satisfied;
    \item product completion time is the finish time of its last required operation.
\end{enumerate}

Let \(C_p\) be the completion time of product \(p\). Product tardiness is
\[
T_p=\max(0,C_p-d_p),
\]
and the primary objective is
\[
\min T^{tot}=\sum_{p\in\mathcal{P}}T_p.
\]
Among solutions with equal total tardiness, the solution with smaller total sample usage
\[
S^{tot}=\sum_{p\in\mathcal{P}}s_p
\]
is preferred. The evaluation tuple is therefore \(Eval=(T^{tot},S^{tot})\) under lexicographic ordering.

\begin{table}[H]
\centering
\caption{Notation used in the proposed TCSSOP framework}
\label{tab:notation_short}
\small
\begin{tabular}{|l|l|}
\hline
\textbf{Symbol} & \textbf{Definition} \\
\hline
\(\mathcal{P}\) & Set of products/projects, indexed by \(p\) \\ \hline
\(\mathcal{C}\) & Set of chambers, indexed by \(c\) \\ \hline
\(\mathcal{E}\) & Set of feasible environment configurations \\ \hline
\(\mathcal{E}^{H}\) & Set of humidity-controlled environments \\ \hline
\(\mathcal{E}^{V}(S)\) & Set of environments with positive voltage-sensitive workload \\ \hline
\(\mathcal{T}\) & Set of test types, indexed by \(t\) \\ \hline
\(\mathcal{T}_p\) & Required test subset of product \(p\) \\ \hline
\(e(t)\) & Environment requirement of selected test variant \(t\) \\ \hline
\(d_p\) & Due date of product \(p\) \\ \hline
\(C_p\) & Completion time of product \(p\) \\ \hline
\(T_p\) & Tardiness of product \(p\), \(T_p=\max(0,C_p-d_p)\) \\ \hline
\(T^{tot}\) & Total tardiness, \(T^{tot}=\sum_{p\in\mathcal{P}}T_p\) \\ \hline
\(s_p\) & Sample count decision for product \(p\) \\ \hline
\(S\) & Sample vector, \(S=\{s_p\}_{p\in\mathcal{P}}\) \\ \hline
\(s_{\min},s_{\max}\) & Lower and upper bounds for sample counts \\ \hline
\(S^{tot}\) & Total sample usage, \(S^{tot}=\sum_{p\in\mathcal{P}}s_p\) \\ \hline
\(\tau_t\) & Processing time of test \(t\) \\ \hline
\(m_{pt}(S)\) & Effective number of operations generated by test \(t\) of product \(p\) \\ \hline
\(X_c\) & Environment assigned to chamber \(c\) \\ \hline
\(x_{ce}\) & Binary variable: 1 if chamber \(c\) is assigned to environment \(e\) \\ \hline
\(M_c\) & Number of stations in chamber \(c\) \\ \hline
\(M^{tot}\) & Total number of stations, \(M^{tot}=\sum_{c\in\mathcal{C}}M_c\) \\ \hline
\(M^{V,tot}\) & Total station capacity in voltage-capable chambers, \(M^{V,tot}=\sum_{c\in\mathcal{C}}M_cv_c\) \\ \hline
\(h_c,v_c\) & Humidity and voltage capability indicators of chamber \(c\) \\ \hline
\(\nu_p\) & Indicator: 1 if product \(p\) requires voltage-compatible execution \\ \hline
\(W_e(S)\) & Total workload assigned to environment \(e\) \\ \hline
\(W_e^V(S)\) & Voltage-sensitive workload assigned to environment \(e\) \\ \hline
\(\phi_e,\phi_e^V\) & Total and voltage-sensitive workload shares of environment \(e\) \\ \hline
\(T_e,T_e^V\) & Total and voltage-sensitive target station capacities \\ \hline
\(z_e,z_e^V\) & Capacity-deviation variables in the assignment model \\ \hline
\(\Pi\) & Feasible schedule \\ \hline
\(Eval\) & Lexicographic objective tuple \((T^{tot},S^{tot})\) \\ \hline
\(I_{\max},K_{\max}\) & Outer-iteration and VNS shake-budget limits \\ \hline
\end{tabular}
\end{table}

\section{Solution Methodology}
\label{sec:solution}

\subsection{Overview}
\label{subsec:framework}

The proposed solution framework coordinates three decision layers: chamber-environment assignment, sample-size optimization, and feasible scheduling. The method starts from the fixed three-sample policy. At each outer iteration, chamber settings are selected for the current sample vector, sample counts are improved under the fixed chamber assignment, and the resulting solution is evaluated by a constructive scheduler. If local improvement stagnates, a VNS diversification phase perturbs sample sizes and re-applies local improvement.

\begin{algorithm}[H]
\caption{Integrated TCSSOP heuristic}
\label{alg:integrated_short}
\begin{algorithmic}[1]
\State Initialize \(S^{current}\) with \(s_p=3\) for all products.
\State Set incumbent \((X^{best},S^{best},\Pi^{best},Eval^{best})\leftarrow(\emptyset,\emptyset,\emptyset,(+\infty,+\infty))\).
\For{\(iter=1,\ldots,I_{\max}\)}
    \State Assign chamber environments \(X^{current}\) for \(S^{current}\).
    \State Apply sample-size local search under \(X^{current}\), obtaining \((S^{local},\Pi^{local},Eval^{local})\).
    \If{\(Eval^{local}\) improves \(Eval^{best}\)}
        \State Update incumbent and continue with \(S^{current}\leftarrow S^{local}\).
    \Else
        \State Apply VNS diversification to \(S^{best}\) under \(X^{best}\).
        \If{VNS improves the incumbent}
            \State Update incumbent and continue.
        \Else
            \State Terminate.
        \EndIf
    \EndIf
\EndFor
\State \Return incumbent solution.
\end{algorithmic}
\end{algorithm}

\subsection{Chamber assignment}
\label{subsec:chamber_assignment}

For a fixed sample vector \(S\), chamber environments are determined from environment-specific workload shares. Let
\[
W^{tot}(S)=\sum_{e\in\mathcal{E}}W_e(S),
\qquad
W^{V,tot}(S)=\sum_{e\in\mathcal{E}}W_e^V(S).
\]
The total workload share is
\[
\phi_e=\frac{W_e(S)}{W^{tot}(S)}.
\]
For voltage-sensitive workload, if \(W^{V,tot}(S)>0\),
\[
\phi_e^V=\frac{W_e^V(S)}{W^{V,tot}(S)}.
\]
Otherwise, \(\phi_e^V=0\) for all \(e\in\mathcal{E}\).

Let
\[
M^{tot}=\sum_{c\in\mathcal{C}}M_c,
\qquad
M^{V,tot}=\sum_{c\in\mathcal{C}}M_cv_c
\]
denote the total station capacity and the station capacity located in voltage-capable chambers, respectively. Total workload targets are scaled by \(M^{tot}\), whereas voltage-sensitive workload targets are scaled only by \(M^{V,tot}\), since voltage-sensitive products can be processed only in voltage-capable chambers. Therefore,
\[
T_e=\phi_e M^{tot},
\qquad
T_e^V=\phi_e^V M^{V,tot}.
\]
Let
\[
\mathcal{E}^V(S)=\{e\in\mathcal{E}: W_e^V(S)>0\}
\]
be the set of environments with positive voltage-sensitive workload.

The assignment model uses proportional workload targets rather than absolute horizon-based capacity requirements. Let \(x_{ce}=1\) if chamber \(c\) is assigned to environment \(e\), and let \(\mathcal{E}^H\) be the set of humidity-controlled environments. The exact assignment model minimizes absolute deviations from workload-based target capacities:
\begin{align}
\min\quad & \sum_{e\in\mathcal{E}}z_e+\sum_{e\in\mathcal{E}}z_e^V \\
\text{s.t.}\quad
& \sum_{e\in\mathcal{E}}x_{ce}=1 && \forall c\in\mathcal{C},\\
& \sum_{c\in\mathcal{C}}x_{ce}\ge 1 && \forall e\in\mathcal{E},\\
& \sum_{c\in\mathcal{C}}v_cx_{ce}\ge 1 && \forall e\in\mathcal{E}^V(S),\\
& x_{ce}\le h_c && \forall c\in\mathcal{C},\ e\in\mathcal{E}^H,\\
& z_e\ge \sum_{c\in\mathcal{C}}M_cx_{ce}-T_e && \forall e\in\mathcal{E},\\
& z_e\ge T_e-\sum_{c\in\mathcal{C}}M_cx_{ce} && \forall e\in\mathcal{E},\\
& z_e^V\ge \sum_{c\in\mathcal{C}}M_cv_cx_{ce}-T_e^V && \forall e\in\mathcal{E},\\
& z_e^V\ge T_e^V-\sum_{c\in\mathcal{C}}M_cv_cx_{ce} && \forall e\in\mathcal{E},\\
& x_{ce}\in\{0,1\},\quad z_e,z_e^V\ge 0 && \forall c,e.
\end{align}
The resulting assignment is then refined by swap and move neighborhoods. A candidate is retained only if it preserves environment coverage, humidity feasibility, and voltage feasibility. Feasible candidates are evaluated with the lightweight EDD evaluation mode described in Section~\ref{subsec:scheduling}, and a first-improvement rule is used.

\subsection{Sample-size local search}
\label{subsec:sample_optimization}

Given chamber assignment \(X\), local search explores product-level sample changes. For each product \(p\), the neighborhood is
\[
\Delta_p \subseteq \{-2,-1,+1,+2\}
\]
subject to sample bounds. Candidate vectors are screened by the lightweight mode of the same EDD scheduler. If a move improves the lexicographic objective, the move is accepted and the current solution is updated.

Both increasing and decreasing sample counts are considered. Increasing a sample count may improve a tardy product by enabling more downstream parallelism. Decreasing a sample count may reduce congestion caused by pull-down operations and may also reduce the secondary objective when tardiness is unchanged.

A Late--Early Rebalance operator is also used. Products with positive tardiness are paired with products having earliness slack. For each pair, one sample is added to the late product and one sample is removed from the early product, subject to bounds. The paired move is accepted if total tardiness does not worsen and, under equal tardiness, total sample usage does not increase.

\begin{algorithm}[H]
\caption{Sample-size local search}
\label{alg:sample_ls_short}
\begin{algorithmic}[1]
\State \textbf{Input:} Chamber assignment \(X\), sample vector \(S\)
\State Clamp all sample values to \([s_{\min},s_{\max}]\).
\State Construct an initial schedule evaluation \(Eval\).
\Repeat
    \State \(improved\leftarrow false\)
    \For{each product \(p\in\mathcal{P}\)}
        \State Define feasible moves \(\Delta_p\subseteq\{-2,-1,+1,+2\}\).
        \State Screen all \(s_p+\delta\), \(\delta\in\Delta_p\), using lightweight EDD evaluation.
        \If{the best screened move improves the incumbent}
            \State Apply the move and update \(Eval\).
            \State \(improved\leftarrow true\)
        \EndIf
    \EndFor
    \State Apply the Late--Early Rebalance operator.
    \If{the rebalance solution is accepted}
        \State Update \(S\) and \(Eval\); set \(improved\leftarrow true\).
    \EndIf
\Until{\(improved=false\) or the pass limit is reached}
\State \Return \((S,Eval)\)
\end{algorithmic}
\end{algorithm}

\subsection{Scheduling, lightweight evaluation, and feasibility validation}
\label{subsec:scheduling}

For fixed \(X\) and \(S\), schedules are constructed using an earliest-due-date (EDD) ready-set rule. Candidate jobs are generated project-wise from the current stage state, sample availability, and earliest feasible station time. The best ready job of each product enters a global candidate pool, and the next job is selected by due date, with earliest feasible start time as a tie-breaker.

Each selected job is assigned to the earliest compatible chamber station. Compatibility requires matching environment, humidity capability when needed, and voltage capability for voltage-sensitive products. Sample availability is updated after every assignment to prevent overlap on the same physical unit.

The fast schedule-aware scoring step is not an aggregate workload proxy and it does not use a different dispatching rule from the final scheduler. It is a lightweight execution mode of the same EDD-based constructive scheduler. For a fixed chamber assignment \(X\) and sample vector \(S\), the scheduler constructs the same ready-set sequence, applies the same chamber, humidity, voltage, station, sample, and stage-state feasibility rules, and returns the resulting total tardiness value. The computational saving comes from suppressing outputs that are not needed during neighborhood screening: the fast mode does not build the detailed operation-level schedule list or, when project-level lateness information is not needed, the project-result list.

More specifically, the score of a candidate is
\[
Score(X,S)=\sum_{p\in\mathcal{P}}\max(0,C_p(X,S)-d_p),
\]
where \(C_p(X,S)\) is obtained from the same deterministic scheduler used for the final evaluation. One lightweight evaluation mode returns only total tardiness and is used when screening room-assignment moves, single-product sample-size moves, and VNS shake candidates. A second lightweight mode returns total tardiness and project-level completion/lateness values but still skips the operation-level schedule list; this mode is used when the Late--Early Rebalance operator requires project-level lateness and earliness information. Accepted outer-loop solutions are then re-evaluated in the schedule-building mode, which constructs the full operation-level schedule and enables the independent feasibility validation checks.

After construction, the schedule is validated independently. The validation checks required-test multiplicity, chamber-environment compatibility, humidity and voltage feasibility, station non-overlap, sample non-overlap, stage-state logic, and completion-time calculation. This layer is essential because infeasible schedules are operationally unusable regardless of objective value.

\begin{algorithm}[H]
\caption{EDD-based schedule construction}
\label{alg:scheduler_short}
\begin{algorithmic}[1]
\State \textbf{Input:} Chamber assignment \(X\), sample vector \(S\)
\State Initialize station availability, sample availability, product states, and schedule \(\Pi\).
\While{at least one product has an executable operation}
    \State Create an empty candidate pool.
    \For{each product \(p\in\mathcal{P}\)}
        \State Identify the best ready operation of \(p\) using stage state, compatible chambers, and sample availability.
        \If{a ready operation exists}
            \State Add it to the candidate pool.
        \EndIf
    \EndFor
    \State Select the next operation by earliest due date; break ties by earliest feasible start.
    \State Assign the operation to the earliest feasible compatible station and sample.
    \State Update station, sample, and product-stage states.
\EndWhile
\State Compute \(C_p\), \(T_p\), \(T^{tot}\), and \(S^{tot}\).
\State Validate all feasibility conditions.
\State \Return \((\Pi,(T^{tot},S^{tot}))\)
\end{algorithmic}
\end{algorithm}

\subsection{VNS diversification}
\label{subsec:vns}

VNS is activated when the outer loop does not improve the incumbent. The method uses a fixed fallback ratio sequence
\[
R^{VNS}=[0.20,0.30].
\]
For each ratio, up to \(K_{\max}\) shake attempts are performed. A shake selects the corresponding fraction of products uniformly without replacement and alternates between two directions: \textsc{DOWN}, which resets selected samples to \(s_{\min}\), and \textsc{UP}, which resets selected samples to the VNS upper level. The perturbed vector is then improved by the same sample-size local search. The first locally improved candidate that improves the incumbent is accepted.

The VNS procedure starts each ratio phase from the incumbent rather than from the previous failed shake. This prevents the search from drifting too far from the best known solution through a sequence of non-improving perturbations.

\begin{algorithm}[H]
\caption{VNS with fixed fallback ratios}
\label{alg:vns_short}
\begin{algorithmic}[1]
\State \textbf{Input:} Chamber assignment \(X\), incumbent \((S^{inc},\Pi^{inc},Eval^{inc})\)
\State Set \(R^{VNS}=[0.20,0.30]\).
\For{each ratio \(r\in R^{VNS}\)}
    \For{\(k=1,\ldots,K_{\max}\)}
        \State \(S^{cand}\leftarrow S^{inc}\).
        \If{\(k\) is odd}
            \State Reset a uniformly selected \(r\)-fraction of eligible products to \(s_{\min}\).
        \Else
            \State Reset a uniformly selected \(r\)-fraction of eligible products to \(s_{\max VNS}\).
        \EndIf
        \State Apply sample-size local search to \(S^{cand}\) under \(X\).
        \If{the locally improved candidate improves \(Eval^{inc}\)}
            \State \Return the candidate and set \(improved=true\).
        \EndIf
    \EndFor
\EndFor
\State \Return the incumbent and set \(improved=false\).
\end{algorithmic}
\end{algorithm}

\section{Computational Study}
\label{sec:computational_study}

\subsection{Experimental design}
\label{subsec:exp_design}

The benchmark set contains 90 scenario combinations generated from \(10\) test matrices, \(3\) due-date scenarios, and \(3\) voltage scenarios. Each test matrix contains 50 product projects. Hence, each benchmark instance contains 50 projects, and the full experimental design is
\[
10 \times 3 \times 3 = 90
\]
instances. The test matrices define the binary required-test vector of each project over the catalog in Table~\ref{tab:test_attributes_short}. The due-date and voltage scenarios are then crossed with these test matrices to separate the effects of portfolio composition, deadline pressure, and voltage-compatibility intensity.

The due-date scenarios differ in both location and dispersion. Scenario~1 is a relatively concentrated deadline setting, with due dates ranging from 42 to 57 days and a median of 50 days; its first and third quartiles are 46 and 54 days, respectively. Scenario~2 introduces a wider spread and earlier deadlines, with due dates ranging from 30 to 69 days, a median of 49.5 days, and first and third quartiles of 40 and 59 days. Scenario~3 represents a more relaxed but still heterogeneous deadline setting, with due dates ranging from 34 to 69 days, a median of 53 days, and first and third quartiles of 43 and 63 days.

The voltage scenarios control the proportion of projects requiring voltage-compatible execution. Scenario~1 corresponds to low voltage intensity, where approximately 15\% of projects require voltage-compatible chambers. Scenario~2 represents medium voltage intensity with a 30\% voltage-requiring share, and Scenario~3 represents high voltage intensity with a 45\% voltage-requiring share. These settings create progressively stronger competition for voltage-capable chambers.

The instances are provided to the Java implementation as a structured input table with one row per project. Each row specifies the instance identifier, project identifier, test-matrix scenario, due-date scenario, voltage scenario, due date, voltage requirement indicator, and the binary required-test columns. The implementation was developed and executed in Java using the Eclipse IDE.

\subsection{Compared methods and parameters}
\label{subsec:methods_params}

Four variants are compared:
\begin{itemize}
    \item \textbf{FS (Fixed Sample):} all products use \(s_p=3\); chamber assignment is performed once and no sample optimization is applied.
    \item \textbf{SO-noOuter:} sample sizes are optimized under the initial chamber assignment, without iterative chamber reassignment.
    \item \textbf{SO-noVNS:} sample optimization and chamber reassignment are iterated, but VNS diversification is disabled.
    \item \textbf{Proposed Framework:} sample optimization, chamber reassignment, and VNS diversification are all enabled.
\end{itemize}

All algorithmic variants were implemented in Java and executed from the Eclipse IDE under the same parameter settings. The chamber-environment assignment model was solved by an in-house exact branch-and-bound routine over chamber types, and the resulting chamber assignments were evaluated by the constructive scheduler described in Section~\ref{subsec:scheduling}. No instance-specific parameter tuning was applied.

All parameters are fixed across instances. The initial sample count is \(s_p^{(0)}=3\), \(s_{\min}=3\), \(s_{\max}=11\), and the VNS \textsc{UP} level is \(s_{\max VNS}=8\). The sample neighborhood is \(\{-2,-1,+1,+2\}\), the VNS fallback ratios are \([0.20,0.30]\), \(I_{\max}=100\), and \(K_{\max}=200\) for the selected configuration. The objective is evaluated lexicographically as \((T^{tot},S^{tot})\).

\begin{table}[H]
\centering
\caption{Parameter settings used in the computational study}
\label{tab:param_settings_short}
\small
\begin{tabular}{lll}
\hline
\textbf{Parameter} & \textbf{Description} & \textbf{Value} \\
\hline
\(s_p^{(0)}\) & Initial sample policy per product & 3 \\
\(s_{\min}\) & Minimum feasible sample count & 3 \\
\(s_{\max VNS}\) & Maximum sample count for VNS UP move & 8 \\
\(s_{\max}\) & Maximum feasible sample count & 11 \\
\(\Delta_p\) & Sample-size neighborhood & \(\{-2,-1,+1,+2\}\) \\
\(R_{\mathrm{VNS}}\) & VNS fallback perturbation ratios & \([0.20,0.30]\) \\
\(I_{\max}\) & Maximum outer-loop iterations & 100 \\
\(K_{\max}\) & Maximum shakes per VNS ratio phase & 200 \\
\hline
\end{tabular}
\end{table}

\subsection{Performance measures}
\label{subsec:metrics}

For method \(m\) and instance \(i\), the percentage gap is computed from total tardiness:
\[
Gap_{im}=
\frac{T^{tot}_{im}-T^{tot,best}_{i}}
{T^{tot,best}_{i}}\times 100,
\]
where \(T^{tot,best}_{i}\) is the best total tardiness observed for that instance in the relevant comparison. In the reported experiments, all denominators are positive. Best-instance counts are based on the lexicographic objective, so total sample usage breaks ties when total tardiness is equal.

\subsection{Results}
\label{subsec:results}

Table~\ref{tab:vns_ratio_comparison_short} compares VNS perturbation policies with \(K_{\max}=100\). The \(0.20\)--\(0.30\) fallback policy produces the lowest average gap and standard deviation among these variants, showing that moderate perturbation followed by stronger fallback is more robust than a single fixed ratio.

\begin{table}[H]
\centering
\caption{Comparison of VNS perturbation policies with \(K_{\max}=100\).}
\label{tab:vns_ratio_comparison_short}
\small
\begin{tabular}{lccccc}
\hline
\textbf{Measure}
&
\shortstack{\textbf{0.10}\\\textbf{VNS (100)}}
&
\shortstack{\textbf{0.20}\\\textbf{VNS (100)}}
&
\shortstack{\textbf{0.30}\\\textbf{VNS (100)}}
&
\shortstack{\textbf{0.10--0.20}\\\textbf{--0.30}\\\textbf{VNS (100)}}
&
\shortstack{\textbf{0.20--0.30}\\\textbf{VNS}\\\textbf{(100)}} \\
\hline
Average gap (\%) & 29.62 & 18.11 & 23.44 & 22.06 & 14.93 \\
Std. dev. (\%) & 44.09 & 31.52 & 30.10 & 33.25 & 27.35 \\
\shortstack[l]{Best-instance\\count} & 25 & 27 & 18 & 38 & 36 \\
\shortstack[l]{Average runtime\\(min)} & 41.3 & 32.1 & 30.4 & 66.0 & 49.7 \\
\hline
\end{tabular}
\end{table}

Table~\ref{tab:vns_budget_comparison_short} evaluates the shake budget for the selected \(0.20\)--\(0.30\) policy. Increasing \(K_{\max}\) improves both average performance and robustness. With \(K_{\max}=200\), the average gap is only \(0.31\%\), the standard deviation is \(1.65\%\), and the method obtains the best result in 82 of 90 scenarios within this comparison.

\begin{table}[H]
\centering
\caption{Effect of shake budget for the \(0.20\)--\(0.30\) VNS policy.}
\label{tab:vns_budget_comparison_short}
\begin{tabular}{lccc}
\hline
\textbf{Measure}
& \shortstack{\textbf{0.20--0.30}\\\textbf{VNS (50)}}
& \shortstack{\textbf{0.20--0.30}\\\textbf{VNS (100)}}
& \shortstack{\textbf{0.20--0.30}\\\textbf{VNS (200)}} \\
\hline
Average gap (\%) & 13.79 & 5.65 & 0.31 \\
Std. dev. (\%) & 29.07 & 20.06 & 1.65 \\
\shortstack[l]{Best-instance\\count} & 18 & 47 & 82 \\
\shortstack[l]{Average runtime\\(min)} & 32.6 & 49.7 & 85.1 \\
\hline
\end{tabular}
\end{table}

The selected proposed configuration is then compared with the main benchmark variants in Table~\ref{tab:main_benchmark_comparison_short}. It obtains the best solution in all 90 scenarios. The fixed-sample policy has the largest average gap, confirming that a uniform sample rule is not robust under heterogeneous due dates, environmental requirements, and chamber-capability constraints. The no-outer and no-VNS variants also remain far from the selected configuration, showing that both assignment--scheduling feedback and diversification are necessary.

\begin{table}[H]
\centering
\caption{Overall comparison of the selected proposed configuration and benchmark variants.}
\label{tab:main_benchmark_comparison_short}
\small
\begin{tabular}{lcccc}
\hline
\textbf{Method}
& \shortstack{\textbf{Average}\\\textbf{gap (\%)}}
& \shortstack{\textbf{Std.}\\\textbf{dev. (\%)}}
& \shortstack{\textbf{Best-instance}\\\textbf{count}}
& \shortstack{\textbf{Average}\\\textbf{runtime (min)}} \\
\hline
Fixed sample & 361.93 & 198.43 & 0 & 0.1 \\
\shortstack[l]{No outer\\reconfiguration} & 140.07 & 90.32 & 0 & 17.1 \\
\shortstack[l]{No VNS\\diversification} & 178.67 & 113.52 & 0 & 0.5 \\
\shortstack[l]{Selected proposed\\configuration} & 0.00 & 0.00 & 90 & 85.1 \\
\hline
\end{tabular}
\end{table}

Table~\ref{tab:fixed_selected_sample_comparison_short} provides a direct comparison between the fixed policy and the selected configuration. The selected configuration uses more samples on average, but it allocates them selectively to products for which additional parallelism reduces tardiness. The goal is therefore not uniform prototype expansion, but coordinated sample allocation and chamber reconfiguration.

\begin{table}[H]
\centering
\caption{Direct comparison between fixed samples and the selected proposed configuration.}
\label{tab:fixed_selected_sample_comparison_short}
\begin{tabular}{lcc}
\hline
\textbf{Measure}
& \textbf{Fixed sample}
& \shortstack{\textbf{Selected proposed}\\\textbf{configuration}} \\
\hline
Average gap (\%) & 361.93 & 0.00 \\
Std. dev. (\%) & 198.43 & 0.00 \\
\shortstack[l]{Best-instance\\count} & 0 & 90 \\
\shortstack[l]{Average runtime\\(min)} & 0.1 & 85.1 \\
\shortstack[l]{Average sample\\count} & 150 & 380 \\
\hline
\end{tabular}
\end{table}

The sample increase observed in the selected configuration should be interpreted within the objective definition. Total tardiness is the primary tactical objective, and total sample usage is used as a secondary lexicographic criterion. Therefore, the method accepts additional samples only when they improve total tardiness or when equal-tardiness alternatives can be distinguished by lower sample usage. In settings where prototype cost or sample availability is a binding managerial constraint, the framework can be extended by adding an explicit sample budget \(\sum_p s_p\le B\) or by using a weighted multi-objective formulation.

\subsection{Managerial implications}
\label{subsec:managerial}

The results provide three managerial insights. First, sample counts should be optimized rather than fixed. The value of an additional sample depends on due-date urgency, downstream parallelization potential, and the compatible chamber set. Second, sample optimization and chamber assignment should be updated together because changing sample counts shifts environment-specific and voltage-sensitive workload. Third, diversification is important in tactical planning: local sample changes alone may not escape a poor workload structure, whereas controlled VNS perturbations can reveal substantially better sample allocations.

The results also suggest that capacity shortages should not automatically be addressed by uniformly increasing prototype counts. Additional samples create both benefits and costs. They increase pull-down workload, but they can also shorten downstream completion times by enabling parallel tests. The useful managerial question is therefore not whether more samples are always better, but where additional samples create the largest reduction in tardiness.

Finally, feasibility validation should be treated as a core part of decision support. A schedule that appears efficient at an aggregate level may fail because of sample overlap, station overlap, or chamber compatibility. The proposed framework separates lightweight evaluation from exact schedule construction and validation, allowing the search to remain efficient while ensuring that reported schedules are operationally meaningful.

\section{Conclusion}
\label{sec:conc}

This study introduced TCSSOP, an integrated test-chamber scheduling problem with decision-dependent workload generation. The problem arises in refrigerator R\&D testing, where heterogeneous chambers, environmental compatibility, voltage and humidity restrictions, stage-state logic, due dates, and physical sample availability interact.

The proposed framework jointly optimizes chamber settings, sample sizes, and feasible schedules. Chamber environments are assigned using workload-balancing logic; sample sizes are improved by local search and Late--Early Rebalance; schedules are constructed by an EDD-based feasible scheduler; and VNS diversification perturbs sample decisions when local improvement stagnates. Feasibility is enforced through explicit checks for required tests, chamber compatibility, humidity and voltage eligibility, station and sample non-overlap, and stage-state correctness.

Computational experiments on 90 benchmark scenarios show that the selected proposed configuration with \(0.20\)--\(0.30\) VNS fallback and \(K_{\max}=200\) consistently outperforms fixed-sample and ablated variants. The fixed three-sample policy is fast but produces substantially higher tardiness. The ablation results show that both iterative chamber reconfiguration and VNS diversification materially improve solution quality.

The main practical conclusion is that prototype sample allocation should be treated as a tactical planning variable. Samples should be increased selectively where additional parallelism reduces tardiness, and chamber configurations should be reassessed after sample decisions change. Future work may incorporate uncertain test durations, test failures and retesting, chamber setup times, explicit prototype-cost budgets, energy and labor objectives, and rolling-horizon project arrivals.

\bibliographystyle{plainnat}
\bibliography{References}

\end{document}
